%% kapa: exam=1 round=37 year=2026 period=1 form=A subject=civil_law subject_ko=민법 seq=35 number=35
%% source: Q-Net 감정평가사 기출문제 (cst00309) artlSeq=5258924 file=raw/1008/5258924/2026년 제37회 감정평가사 1차 1교시.pdf
%% conversion: kordoc 4.24.51 -> tools/gen1.py (auto); review=proofread-sonnet; flags=none
\begin{question}{35}
甲 소유의 X토지를 임차하여 그 위에 Y건물을 소유하고 있는 乙은 丙과 Y건물에 관한 전세권 설정계약을 체결하였다. 이에 관한 설명으로 옳지 않은 것은? (다툼이 있으면 판례에 따름)
\begin{choices}
  \item 丙 명의의 전세권의 효력은 Y건물의 소유를 목적으로 한 X토지의 임차권에는 미치지 않는다.
  \item Y건물의 일부가 丙에게 책임 있는 사유로 멸실된 경우, 乙은 丙의 전세권이 소멸된 후 전세금으로써 손해의 배상에 충당할 수 있다.
  \item 丙 명의의 전세권 설정과 동시에 Y건물이 丙에게 인도되지 않았더라도 장차 丙이 Y건물을 사용ㆍ수익하는 것을 완전히 배제하지 않았다면, 그와 같은 丙 명의의 전세권의 효력을 부인할 수 없다.
  \item 丙 명의의 전세권이 설정된 이후에 乙이 丁에게 Y건물의 소유권을 이전한 경우, 丁은 丙에 대하여 전세권 설정자의 지위에서 전세금 반환의무를 부담하게 된다.
  \item 丙 명의의 전세권이 Y건물의 일부에 설정된 경우, 丙은 전세권의 목적물이 아닌 나머지 부분에 대해서는 경매를 신청할 수 없다.
\end{choices}
\end{question}
\answer{1}
